\section{定位与设计原则}
\label{sec:positioning}

\subsection{与既有 TLRA 脉络的连续性}

既有 TLRA 线路确立了三个仍居中心地位的原则。第一，幻觉缓解应作用于\textbf{生成过程之中}，而非仅依赖事后过滤。第二，控制的基本单元应是\textbf{候选局部（candidate-local）}的，因为幻觉往往出现在具体词法替代层面而非全局序列质量。第三，任何实用的解码时干预必须保持语言结构，而非无差别压制视觉薄弱 token。

UniGround 在改变可学习组件的位置与角色的同时保留上述原则。它被表述为解码时接地控制接口，而非关于证据应如何在宿主 MLLM 内部重路由的理论。目标不是重新解释宿主模型的内部推理轨迹，而是在严格通用性约束下，附着可迁移的外部控制器，以评估并调节候选片段。

\subsection{为何内部词表对齐不是``通用''的}

当模块的输入或输出是宿主特有量时，学习模块在严格意义上不具通用性。更精确地说，通用性在以下任一成立时即被破坏：
\begin{enumerate}
  \item 模块消费其语义依赖宿主架构的隐状态；或
  \item 模块输出被直接解释于宿主模型词表空间中的向量。
\end{enumerate}

该点可直接形式化。设 $m$ 为宿主 MLLM，其隐状态空间为 $\mathcal{H}_m$，词表为 $\mathcal{V}_m$，反嵌入映射为 $U_m$。假设学习控制器 $f_\theta$ 消费 $h_t \in \mathcal{H}_m$，并在宿主词表空间 $\mathbb{R}^{|\mathcal{V}_m|}$ 中输出修正。若不存在模型无关的同构，能同时对齐两宿主 $m$ 与 $m'$ 的 $\mathcal{H}_m,\mathcal{H}_{m'},\mathcal{V}_m,\mathcal{V}_{m'}$，则同一参数 $\theta$ 无法在两模型上保持语义等价。因而该学习映射按构造即为宿主耦合。

\texttt{Phi\_calib} 恰好实例化这一宿主耦合模式：它读取模型原生的多模态隐状态并投影到当前模型的词表几何。其行为因此依赖隐状态来源、词表结构、分词器约定与 \texttt{lm\_head} 组织。它作为内部对照仍有信息量，但不能支撑严格的通用插件主张。

\subsection{通过标准化观测—动作空间实现通用性}

要使通用性可辩护，学习插件应通过模型无关的观测与动作空间定义：
\begin{itemize}
  \item \textbf{观测空间：} 像素、冻结图像嵌入、冻结文本嵌入、解码字符串，以及可选的冻结区域提议。
  \item \textbf{动作空间：} 片段级的视觉支持、矛盾与弃权估计。
\end{itemize}

在此设计下，学习插件从不消费模型原生隐状态，不读取 \texttt{lm\_head.weight}，也不依赖学习式的模型专用桥接。通用性因而在表示与附着两个层面共同被强制执行。方法之所以``通用''，并非因为每个组件跨模型完全相同，而是因为所有模型相关变化被限制在\textbf{确定性的、非学习的}附着逻辑中。
